\documentclass[10pt,a4paper]{amsart}
\usepackage[margin=19mm]{geometry}
\usepackage{amsmath,amssymb,mathtools}
\usepackage[scr=rsfs,cal=euler]{mathalfa}
\usepackage{xfrac}
\usepackage[unicode,hidelinks,bookmarks=false]{hyperref}
\newcommand{\ie}{i.e.\ }
\newcommand{\cf}{cf.\ }
\newcommand{\resp}{respectively\ }
\newcommand{\Subsection}{\subsection}
\providecommand{\nobreakdash}{\nobreak\hbox{-}}
\setlength{\emergencystretch}{2em}
\multlinegap0pt
\usepackage{tikz}
\usetikzlibrary{cd}
\usepackage{xcolor}
\usepackage{tikz}
\usepackage{enumerate}
\usepackage[shortlabels]{enumitem}
\usepackage[normalem]{ulem}
\usepackage{amssymb}
\usepackage{mathtools}
\usepackage{float}
\newtheorem{proposition}{Proposition}
\newcommand{\VecSp}{\mathbf{Vec}}
\newcommand{\im}{\operatorname{im}}
\newcommand{\imapx}[1]{j_{#1\ast}}
\newcommand\michal[1]{\textcolor{ForestGreen}{[ML: #1]}}
\renewcommand{\phi}{\varphi}
\title{Connection Matrix: Efficient Computation and Persistence}
\author{Tamal K. Dey, Andrew Haas, Micha\l{} Lipi\'nski}
\date{Source: arXiv:2608.06507v1 (CC BY 4.0)}
\begin{document}
\maketitle

\noindent\emph{Source and licence.} Connection Matrix: Efficient Computation and Persistence. Version arXiv:2608.06507v1 (CC BY 4.0), distributed by arXiv under the Creative Commons Attribution 4.0 International licence, http://creativecommons.org/licenses/by/4.0/, as recorded in the arXiv licence metadata for this exact version and shown on the abstract page https://arxiv.org/abs/2608.06507v1. Full licence notice: \texttt{licenses/arxiv-2608.06507-CC-BY-4.0.txt}.\par\smallskip

\noindent\textit{Editorial scope.} Counted unit: the article's proposition prop:CMmodule-simplified (source0.tex lines 1412--1418). The article's complete original proof of it is delivered with it (source0.tex lines 1419--1470, one \texttt{proof} environment of 52 lines). No mathematics is written, repaired or re-derived: the statement and the proof are the article's own lines, in the article's own notation, and no retained line is substituted or re-flowed.  Retained uncounted as the source-local context the proof needs: lines 1263--1267; lines 1269--1279; lines 1369--1375.  Nothing was omitted from any retained span, so the package carries no omission record.  No retained line contains a citation, so the wrapper carries no bibliography and no bibliography entry.  The retained lines contain the article's own diagram code, which the wrapper typesets with the standard tikz packages; no image file is delivered and the package declares no asset dependency.  Editorial formatting only, in addition: an AMS \texttt{amsart} preamble that restates the article's own declarations and macros used by the retained lines, namely the environments \texttt{proposition} on separate counters, together with the standard packages those lines need.  The retained environments print in this document as Proposition 1 in order of appearance, whereas the article prints the same items with the numbering of its own full document.  The complete native source of the article as distributed in the arXiv e-print for this exact version is bundled unchanged as \texttt{sources/207004/source0.tex}.
\begin{proposition}\label{prop:structure-of-i-j-ast}
    Restrictions $i_\ast|_{W_1}:W_1\rightarrow X_1$ and $j_\ast|_{X_2}:X_2\rightarrow V_2$ are isomorphisms.  
    % $i_\ast|_{W_1}$ and $j_\ast|_{X_2}$ are isomorphisms.  
    Moreover, $i_\ast|_{W_2}=0$ and $j_\ast|_{X_1}=0$.
\end{proposition}

\begin{equation}\label{eq:lefschetz_ar_split_joint}
    \begin{tikzcd}[row sep=1.5em, column sep=1.5em]
        & & & H(B_2)\ar[r,equal]& V_1\oplus V_2\\
        & & & H(D_2, E_2)\arrow[u, "\varphi_2"]\ar[r,equal]& V_1'\oplus V_2'\\
        X_1\oplus X_2\ar[r, equal] & H(B) \arrow[rruu, bend left, "j_\ast"]\arrow[r,leftarrow, "\psi"]
            & H(D,E)\arrow[ru, "j_\ast'"] & \ar[l, equal]X_1'\oplus X_2' \\
        & & & H(D_1, E_1)\arrow[d, "\varphi_{1}"]\arrow[lu, "i_\ast'"]\ar[r,equal]& W_1'\oplus W_2'
            \\ %W_1\oplus W_2 \cong\im \imapx{1}\oplus\ker\imapx{1}\ar[l, equal]\\
        & & & H(B_1)\arrow[lluu, bend left, "i_\ast"] \ar[r,equal]& W_1\oplus W_2
    \end{tikzcd}
\end{equation}

\begin{proposition}
    Let $M:\mathbb X\rightarrow \VecSp$ be a persistence module on a finite
    poset $\mathbb X$ and let $\{\mathbb{I}_\alpha\}_{\alpha\in \Lambda}$ be a set of submodules of $M$ with
    $\bigoplus_{\alpha\in\Lambda}\mathbb{I}_\alpha(x)=M(x)$  $\forall x\in \mathbb{X}$.
    Then, $M=\bigoplus_{\alpha\in\Lambda} {\mathbb I}_\alpha$.
    \label{prop:composition}
\end{proposition}

\begin{proposition}
    Let $\overline\Pi$ be the set of paths in $\overline{\mathbb{P}}$ so that %${\mathcal G}=\bigoplus_{\pi\in \Pi} \mathbb{I}_{\pi}$, and 
    $\overline{\mathcal G}=\bigoplus_{\pi\in \overline{\Pi}} \bar{\mathbb{I}}_{\pi}$.
    Then, ${\mathcal G}=\bigoplus_{\pi\in \Pi} \mathbb{I}_{\pi}$
    for a set of paths $\Pi$ in $\mathbb{P}$ where $\Pi=\overline{\Pi}$.
    \label{prop:CMmodule-simplified}
\end{proposition}

\begin{proof}
    First, we recognize that both $\mathbb{P}$ and $\mathbb{\overline{P}}$ have the same 
    set of objects (points) and a one-to-one correspondence between their relations in their
    Hasse diagrams (transition diagrams). The only difference is that certain relations in these
    Hasse diagrams have opposite directions, namely, an inclusion $(D_{p,t},E_{p,t})\hookrightarrow (D_{q,t\pm 1},E_{q,\pm 1})$ in index pairs may correspond to an inclusion $B_{p,t}\hookleftarrow B_{q,t\pm 1}$ in blocks.

    Consider any summand $\bar{\mathbb{I}}_\pi$ in the decomposition $\overline{\mathcal G}=\bigoplus_{\pi\in \overline{\Pi}} \bar{\mathbb{I}}_{\pi}$. Let $(p,t)\rightarrow (q,t\pm 1)$ be
    any edge in the Hasse diagram of $\overline{\mathbb{P}}$. 
    % We have $\bar v_p\stackrel{\bar\iota_\ast}{\rightarrow}\bar v_q$
    We have $\bar v_p\stackrel{\bar\kappa_\ast}{\rightarrow}\bar v_q$
    where $\bar v_p=\bar{\mathbb{I}}_\pi(p,t)$ and $\bar v_q=\bar{\mathbb{I}}_\pi(q,t\pm 1)$ are zero or one dimensional vector spaces with the structural map $\bar\kappa_\ast$ in $\overline{\mathcal G}$
        % \michal{I'm not sure about $\iota$ notation here, because we use $\iota$ on the level of blocks, here we are still at the level of index pairs. How about $\bar{i}_\ast$?} 
        induced by inclusions in index pairs. 
    Now, we will show that that the summand $\bar{\mathbb{I}}_\pi$ is consistent with the module $\mathcal G$ defined on~$\mathbb{P}$. 
    % Let us look at the module $\mathcal G$ defined on~$\mathbb{P}$. 
    Recall the definition of the structural maps $h_{p,q}$ for $\mathcal G$. We have two cases.

    % \michal{how do we know that $\bar{\iota}_\ast$ agrees with $i_\ast'$ on $\bar{v}_p$? It is because we assume that $\mathbb{I}$ is a summand, right?}
    \textbf{No reversal.} $(p,t)\rightarrow (q,t\pm 1)$ is
    a relation in $\mathbb P$: by definition 
        $h_{p,q}= i_\ast=\psi\circ i_\ast'\circ\phi_1^{-1}$
        % $h_{p,q}= i_\ast=\psi\circ\bar\iota_*\circ\phi_1^{-1}$
    where $\psi$ and $\phi_1$ are isomorphisms (see Diagram~\ref{eq:lefschetz_ar_split_joint}).
        % , $\bar\iota_*=i_\ast'$). 
    % \michal{
    Since $\bar{\kappa}_\ast=i_\ast'$ we have $h_{p,q}(\phi_1(\bar v_p))=\psi(i_\ast'(\bar v_p))=\psi(\bar v_q)$.
    % } 
    Writing
    $\phi_1(\bar v_p)=v_p\in W$ and $\psi(\bar v_q)=v_q\in  X$, we get $h_{p,q}(v_p)=v_q$.
    % Then, $h_{p,q}(\phi_1(\bar v_p))=\psi(\bar\iota_*(\bar v_p))=\psi(\bar v_q)$. 
    % Writing
    % $\phi_1(\bar v_p)=v_p\in W$ and $\psi(\bar v_q)=v_q\in  X$, we get $h_{p,q}(v_p)=v_q$.

    \textbf{A reversal.} $(p,t)\leftarrow (q,t\pm 1)$ is a relation in $\mathbb P$: 
        If $\bar v_p\in X_1'$, we have $j'_\ast(\bar v_p)=0$ (by Proposition~\ref{prop:structure-of-i-j-ast}), 
            and also $v_p\not\in\im h_{q,p}=\im f$ by construction of $f$. 
            % but also $h_{q,p}(v_p)=0$ by construction of $f$. 
    If $\bar v_p\in X_2'$, then by definition,
        $h_{q,p}(\phi_2(\bar v_q))=\psi(\bar v_p)$ because $j_\ast'(\bar v_p)=\bar v_q$ which agrees with $\bar{\kappa}_\ast$.
        Then, we have $h_{q,p}(v_q)=v_p$.
        
    % \textbf{A reversal.} $(p,t)\leftarrow (q,t\pm 1)$ is a relation in $\mathbb P$: If $\bar v_p\in \ker \bar\iota_*=\ker j_\ast'$, we have $v_p\not\in \im  h_{q,p}$ by construction of $h$. 
    % If $\bar\iota_*(\bar v_p)=\bar v_q\neq 0$, then by definition,
    % $h_{q,p}(\phi_2(\bar v_q))=\psi(\bar v_p)$ because $j_\ast'(\bar v_p)=\bar v_q$.
    % Then, we have $h_{q,p}(v_q)=v_p$.

    It follows that for each summand $\bar{\mathbb{I}}_\pi$ in $\bar{\mathcal G}$, we have a submodule $\mathbb{I}_\pi$ of $\mathcal G$ with the same support. Furthermore, because of the isomorphisms $\psi,\phi_1,\phi_2$,
    $\bigoplus_{\pi\in \Pi} (v_p=\mathbb I_\pi(p,t))=\bigoplus_{\pi\in \Pi}(\bar v_p= \overline{\mathbb{I}}_\pi(p,t))=\overline{\mathcal G}(p,t)=\mathcal{G}(p,t)$.
    Therefore, by Proposition~\ref{prop:composition}, $\mathcal G=\bigoplus_{\pi\in \Pi} \mathbb I_\pi$. By the uniqueness of decomposition up to isomorphism,
    any other decomposition $\mathcal G=\bigoplus_{\pi\in \Pi'}\mathbb I_\pi$ must have
    $\Pi=\Pi'$. \qed
\end{proof}

\end{document}
